\scanpage{14}
\mspage{\handwritten{28}}

1) $\mathrm B$ est formellement simple sur $\mathrm A$

\deleted{\unsure{3}}2) $\mathrm B$ satisfait la condition $\mathrm P$\quad 3) $\mathrm B$ est plat sur $\mathrm A$\quad 3 bis) Pour tout entier $n\handwritten{\geq}0$, $\mathrm B_n=\mathrm B/\mathrm I^{n+1}\mathrm B$ est plat sur $\mathrm A_n=\mathrm A/\mathrm I^{n+1}$. \deleted{4)}

Il reste seulement à montrer l'équivalence de 3) et 3 bis), qui est contenue en effet (moyennant les hypothèses noethériennes faites) dans la platitude supérieure.

\deleted{Le résultat précédent} Il conviendrait de noter aussi\inserted{ dès après 3.7.} que ``algèbre de Cohen'' est une notion stable par extension de la base via produits tensoriels complétés: c'est une trivialité.

Le résultat \deleted{\unsure{3}}\handwritten{4}.8. montre l'intérêt qu'il y a, $\mathrm A,\mathrm I$ et $\mathrm B_{\unsure{0|o}}$ étant donnés, de disposer d'un \uline{théorème d'existence} pour une algèbre de Cohen $\mathrm B$ sur $\mathrm A$ se réduisant suivant $\mathrm B_0$. On notera que la construction habituelle de proche en proche nous ramène \deleted{au cas où $\mathrm A$ est discr}(remplaçant $\mathrm A$ par des $\mathrm A_i$ et $\mathrm I$ par des puissances) au cas où $\mathrm I$ est de carré nul,(donc où $\mathrm A$ est discret). On obtient alors un problème de nature essentiellement infinitésimale, dont nous ignorons s'il existe toujours une solution, faute de savoir caractériser la structure des algèbres de Cohen $\mathrm B_0$ sur l'anneau discret $\mathrm A_0$. Cependant, nous obtenons \inserted{\handwritten{(\deleted{\unsure{3}}\unsure{4|3}.10.)}} un théorème d'existence dans un cas particulier intéressant, en utilisant les résultats du N\textsuperscript{\uline{o}}2\deleted{\unsure{,}}. Notons \inserted{d'abord}

\deleted{\uline{Théorème} 3.9. Soit $\mathrm A$ \deleted{$\mathrm B$ un homomorphisme} un anneau local noethérien, \deleted{et de corps} résiduel $k$, et soit $\mathrm B_0$ une $k$-algèbre, \deleted{qui soit \unsure{u|n} satisfasse} \inserted{\deleted{qui} soit} \deleted{aux conditions suivantes: $\mathrm B_0$ est} un anneau local noethérien, \deleted{et complet,} \inserted{de} \deleted{son corps} résiduel $\mathrm K$ \deleted{est de degré de multiplicité radicielle finie sur $k$ (ce qui est le cas} une extension séparable de $k$, $\mathrm B_0$ est (pour sa topolo-}

\uline{Corollaire} \deleted{3}\handwritten{4}.9. Soit $\mathrm A\handwritten{\to}\mathrm B$ un homomorphisme local d'anneaux locaux noethériens, \deleted{tel que l'extension résiduelle $k(\mathrm B)/k(\mathrm A)$ soit à multiplicité}
% archive sheet number: 13
